% page 110 of the facsimile (image p-109 of the scan)
% block 162.6 x 237.7 mm ; left margin 8.0 mm
\pgc{-12.700mm}{0.000mm}{
\l{\cel{12}lO2}
\l{}
\l{\sou{demisionar}. (\sou{netrans.})\cel{2}Retretar volunte de ofico. - DFIS.}
\l{}
\l{\sou{demiurgo}. (En\cel{1}\sou{la filozofio olima)}.\cel{2}Deo, konsiderata ne kom}
\l{\cel{3}kreinto, ma kom ordininto, di materio. - DEFIRS.}
\l{}
\l{\sou{demokrata}. Esar partisano di demokratio, e (plu generale) havar}
\l{\cel{3}la mento demokratiala o demokratismala. - DEFIRS.}
\l{}
\l{\sou{demokratio}. La rejimo politikala en qua la populo guvernesas}
\l{\cel{3}da su ipsa (t.e. per sua elektiti) : lua principo esas la}
\l{\cel{3}suvereneso di la majoritato. - DEFIRS.}
\l{}
\l{\sou{demolisar}. (\sou{trans.})\cel{2}Deskonstruktar, per faligar, sucedante,}
\l{\cel{3}omna parti qui kompozas lu. - DEFIS.}
\l{}
\l{\sou{demono}. I. (\sou{Che la Antiqui}).\cel{2}Feo bonfacera o malfacera qua-}
\l{\cel{3}segun la kredo lora - prezide direktis la fato di homo,}
\l{\cel{3}di populo, ed inspiris a lu bonaji o malaji. - II. (\sou{Nun})}
\l{\cel{3}Diablo. - III. (\sou{Metaf.)}\cel{3}Personiguro di defekto, di vicio.}
\l{\cel{3}- DEFIRS.}
\l{}
\l{\sou{demonstrar}. (\sou{trans.}) I. Pruvar, per rezono, la exakteso di pro-}
\l{\cel{3}poziciono. - II. Docar per montrar la kozi, per pozar li}
\l{\cel{3}avan la regardo. - DEFIRS.- DEFIRS.}
\l{}
\l{\sou{demonstrativa}. (\sou{Gram.})\cel{2}Qua indikas persono o kozo. - DEFIRS.}
\l{}
\l{\sou{demotika}. (Skribo-formo) kursiva, vulgara, che la Egiptiani}
\l{\cel{3}antiqua, simpliguro di la skriboformo hieratika (hierogli-}
\l{\cel{3}fi). - DEF.}
\l{}
\l{\sou{denaro}. En Roma (antique), monet-uno ek arjento qua valoris}
\l{\cel{3}cirkume dek \textquotedbl{}as-i\textquotedbl{}. - DEFIRS.}
\l{}
\l{\sou{denaturar}.\cel{2}(\sou{trans.})\cel{2}Igar neapta kom nutrivo, ma lasar sat}
\l{\cel{3}apta pri uzeso en industrio. - DEFIRS.}
\l{}
\l{\sou{dendrito}. Petro sur qua esas videbla desegnuri naturala qui}
\l{\cel{3}reprezentas arboro-rami : ta ipsa rami.}
\l{}
\l{\sou{denigrar}. (\sou{trans.})\cel{2}Inspirar mala opiniono pri ulu (od ulo)}
\l{\cel{3}per aserti ed opinioni severa o nejusta. - F.}
\l{}
\l{\sou{denominatoro.}\cel{2}(\sou{matem.})\cel{2}Ta ek la du termini di fraciono qua}
\l{\cel{3}indikas a quanta parti inter-egala dividesis la unajo. -EFIS.}
\l{}
\l{\sou{densa}. Di qua la maso esas granda kompare a la volumino. -EFIS.}
\l{\cel{3}- Relato di la pezo di korpo (solida o liquida) a la pezo di}
\l{\cel{3}volumino egala de aquo distilita ye 4o C. Relato di la pezo}
\l{\cel{3}di gaso a la pezo di volumino egala di aero sama-faktore.}
\l{}
\l{\sou{densimetro}. Aparato por mezurar nemediate la denseso-grado di}
\l{\cel{3}la korpi (precipue la liquidi). - DEFIRS.}
}
